% SHARD-ID: CA-21-LATTICE-GALILEAN-BOOSTS
% SHARD-TITLE: Lattice Galilean Boosts from Mass-Density First Moments
% SHARD-SUMMARY: Derives the Galilean lattice boost candidate as a first moment of conserved mass or particle-number density, not energy density.
% SHARD-SUMMARY: Uses a discrete continuity equation to derive the current-sum candidate for momentum, with coefficients checked in Julia.
% SHARD-SUMMARY: Separates exact finite-lattice conservation from the stronger continuum claim of Galilean covariance.
% SHARD-KEYWORDS: Galilean lattice boost, mass density, particle number, continuity equation, current, momentum candidate

\section{Lattice Galilean Boosts from Mass-Density First Moments}
\label{sec:lattice-galilean-boosts}

\subsection{Status, Sources, and Dependencies}

\textbf{Status: proposal plus checked local derivation.}  Local sources support
Galilean invariance as a physical assumption in nonrelativistic fermion models
and give examples where conservation laws produce continuity equations.  The
finite open-chain first-moment formula below is a local derivation checked by
the Julia test suite.

Dependencies: CA-18--CA-20 and CONVENTIONS.md (h), (i).

% Source: literature/md/cond-mat_9906453/cond-mat_9906453.md:39 -- interacting fermions of mass m with Galilean invariance in zero magnetic field.
% Source: literature/md/cond-mat_9906453/cond-mat_9906453.md:41 -- center of mass accelerated by a uniform electric field while relative motion is unaffected.
% Source: literature/md/cond-mat_9906453/cond-mat_9906453.md:662 -- current density with 1/m factor in a nonrelativistic model.
% Source: literature/md/cond-mat_9906453/cond-mat_9906453.md:664 -- spin conservation implies a continuity equation as an operator equation.
% Source: literature/md/cond-mat_9906453/cond-mat_9906453.md:682 -- continuity equation and gauge invariance imply transverse response constraints.
% Source: literature/md/cond-mat_9906453/cond-mat_9906453.md:716 -- Ward identity from the continuity equation.
% Checked: test/runtests.jl, "first-moment continuity current coefficients".

\subsection{Input Data}

Let \(\Lambda=\{1,\ldots,N\}\) be a finite open chain with positions \(x_j\).
A Galilean lattice candidate needs, in addition to the Hamiltonian, a conserved
mass or particle-number density:
\[
  M=\sum_j \mu_j,
  \qquad
  [H,M]=0.
\]
For identical particles of mass \(m\), a common special case is
\[
  \mu_j=m n_j,
\]
where \(n_j\) is a number density.

The time-zero Galilean boost candidate is the first mass moment
\[
  G^{\mathrm{lat}}(0)=\sum_j x_j\mu_j.
\]
This is the Galilean analogue of \(mX\), not the Lorentzian energy moment
\(\sum_j x_jh_j\).

\subsection{Discrete Continuity Assumption}

Assume a nearest-neighbour open-chain continuity equation with bond currents
\(J_{j+1/2}\):
\[
  i[H,\mu_j]=J_{j-1/2}-J_{j+1/2},
\]
where endpoint currents are either absent or treated as boundary terms.  This
assumption is model data.  It does not follow merely from writing
\(H=\sum_jh_j\).

\subsection{First-Moment Derivation}

Multiply the continuity equation by \(x_j\) and sum:
\[
\begin{aligned}
  i\left[H,\sum_j x_j\mu_j\right]
  &=
    \sum_j x_j(J_{j-1/2}-J_{j+1/2})\\
  &=
    \sum_{j=1}^{N-1}(x_{j+1}-x_j)J_{j+1/2}
    \quad+\quad\text{boundary terms}.
\end{aligned}
\]
For a unit-spaced open chain away from endpoints,
\[
  i[H,G^{\mathrm{lat}}(0)]
  \simeq
  \sum_{j=1}^{N-1}J_{j+1/2}.
\]
This is the lattice momentum candidate associated with the Galilean boost.
For nonuniform positions the coefficients are the spacings
\[
  \Delta x_j=x_{j+1}-x_j,
\]
and the test suite records a coefficient check for this shard.

\subsection{Relation to Center of Mass}

The local paired-fermion source uses Galilean invariance to explain a response
in terms of center-of-mass acceleration while relative motion is unaffected.
That is exactly the intuition this shard turns into compiler data:
\[
  \text{Galilean boost} \quad\leadsto\quad
  \text{center-of-mass / mass-first-moment generator}.
\]
The source does not prove the lattice formula above.  It supports the physical
role of center-of-mass motion under Galilean invariance; the lattice identity is
the checked local derivation.

\subsection{What Is and Is Not Proven}

The finite identity proves only that, given a local continuity equation, the
time derivative of the mass first moment is a current sum.  It does not prove:
\[
\begin{array}{ll}
1.&\text{that the current sum converges to continuum momentum},\\
2.&\text{that the finite lattice has exact Galilei symmetry},\\
3.&\text{that the mass current equals momentum density in every model},\\
4.&\text{that boundary terms vanish without a stated boundary condition}.
\end{array}
\]
Those are acceptance-test obligations for later shards or model-specific
examples.
